\label{sec:beta}
We now consider the problem of inferring borrowing signatures, i.e. a mapping
$\beta : \ty{Const} \rightharpoonup \{\own{}, \bor{}\}^*$, which for every
function should return a list describing each parameter of the function as
either \own{}{}wned or \bor{}{}orrowed.
Borrow annotations can be provided manually by users (which is always safe), but we have two motivations
for inferring them: avoiding the burden of annotations, and making our IR a convenient target
for other systems (e.g., Coq, Idris, and Agda) that do not have borrow annotations.

If a function $f$ takes a parameter $x$ as a borrowed reference, then
at runtime $x$ may be a shared value even when its reference counter
is $1$. Thus, we must never mark $x$ as borrowed if it is used by
a $\kw{let}\ y = \kw{reset}\ x$ instruction. We also assume that
each $\beta(c)$ has the same length as the corresponding parameter
list in $\delta(c)$.

Partially applying constants with borrowed parameters is also
problematic because, in general, we cannot statically assert that the
resulting value will not escape the current function and thus the
scope of borrowed references. Therefore we extend $\delta\reuse$ to
the program $\delta_\beta$ by defining a trivial wrapper
constant $c_\own{} := c$ (we will assume that this name is fresh) for any
such constant $c$, set $\beta(c_\own{}) := \many{\own{}}$, and replace any
occurrence of $\Pap c \many{y}$ with $\Pap {c_\own{}} \many{y}$.
The compiler step given in the next subsection will,
as part of the general transformation,
insert the necessary \kw{inc} and \kw{dec} instructions into
$c_\own{}$ to convert between the two signatures.

Our heuristic is based on the fact that when we mark a parameter as
borrowed, we reduce the number of RC operations needed, but we
also prevent $\kw{reset}$ and $\kw{reuse}$ as well as primitive
operations from reusing memory cells.  Our heuristic collects which
parameters and variables should be owned.  We say a parameter $x$
should be owned if $x$ or one of its projections is used in a
\kw{reset}, or is passed to a function that takes an owned reference.
The latter condition is a heuristic and is not required for
correctness.  We use it because the function taking an owned reference
may try to reuse its memory cell.  A formal definition is given in
\cref{fig:collect_O}. Many refinements are possible, and we
discuss one of them in the next section.
\begin{figure}
\begin{lstlisting}
$\textit{collect}_\own{}$ : $\ty{FnBody}\RC \rightarrow 2^{\ty{Vars}}$
$\textit{collect}_\own{}$(let z = ctor_i $\many{x}$; F) = $\textit{collect}_\own{}(F)$
$\textit{collect}_\own{}$(let z = reset x; F) = $\textit{collect}_\own{}(F) \cup \{x\}$
$\textit{collect}_\own{}$(let z = reuse x in ctor_i $\many{x}$; F) = $\textit{collect}_\own{}(F)$
$\textit{collect}_\own{}$(let z = c $\many{x}$; F) = $\textit{collect}_\own{}(F) \cup \{x_i \in \many{x}\ |\ \beta(c)_i = \own{} \}$
$\textit{collect}_\own{}$(let z = x y; F) = $\textit{collect}_\own{}(F) \cup \{x, y\}$
$\textit{collect}_\own{}$(let z = pap $c_\own{}$ $\many{x}; F$) = $\textit{collect}_\own{}(F) \cup \{\many{x}\}$
$\textit{collect}_\own{}$(let z = proj_i x; F) = $\textit{collect}_\own{}(F) \cup \{x\}$ $\text{if } z \in \textit{collect}_\own{}(F)$
$\textit{collect}_\own{}$(let z = proj_i x; F) = $\textit{collect}_\own{}(F)$ $\text{if } z \not\in \textit{collect}_\own{}(F)$
$\textit{collect}_\own{}$(ret x) = $\emptyset$
$\textit{collect}_\own{}$(case x of $\many{F}$) = $\bigcup_{F_i \in \many{F}} \textit{collect}_\own{}(F_i)$
\end{lstlisting}
\caption{Collecting variables that should not be marked as borrowed\label{fig:collect_O}}
\end{figure}
Note that if a call is recursive, we do not know which parameters are owned, yet. Thus, given
$\delta(c) = \lambda \many{y}.\ b$, we infer the value of $\beta(c)$ by starting
with the approximation $\beta(c) = \bor{}^n$, then we compute $S = \textit{collect}_\own{}(b)$,
update $\beta(c)_i := \own{}$ if $y_i \in S$, and
repeat the process until we reach a fix point and no further updates
are performed on $\beta(c)$. The procedure described here does not
consider mutually recursive definitions, but this is a simple
extension where we process a block of mutually recursive functions
simultaneously.
By applying our heuristic to the $\textit{hasNone}$ function described before, we obtain $\beta(\textit{hasNone}) = \bor{}$. That is, in an application $\textit{hasNone}\ \textit{xs}$, $\textit{xs}$ is taken as a borrowed reference.

%%% Local Variables:
%%% mode: latex
%%% TeX-master: "refcount"
%%% End:
